\documentclass[a4paper,twoside]{article}
\usepackage[T1]{fontenc}
\usepackage[bahasa]{babel}
\usepackage{graphicx}
\usepackage{graphics}
\usepackage{float}
\usepackage[cm]{fullpage}
\pagestyle{myheadings} 
\usepackage{etoolbox}
\usepackage{setspace} 
\usepackage{lipsum} 
\setlength{\headsep}{30pt}
\usepackage[inner=2cm,outer=2.5cm,top=2.5cm,bottom=2cm]{geometry} %margin
% \pagestyle{empty}

\makeatletter
\renewcommand{\@maketitle} {\begin{center} {\LARGE \textbf{ \textsc{\@title}} \par} \bigskip {\large \textbf{\textsc{\@author}} }\end{center} }
\renewcommand{\thispagestyle}[1]{}
\markright{\textbf{\textsc{AIF184001/AIF184002 \textemdash Rencana Kerja Tugas Akhir \textemdash Sem. Ganjil 2026/2027}}}

\newcommand{\HRule}{\rule{\linewidth}{0.4mm}}
\renewcommand{\baselinestretch}{1}
\setlength{\parindent}{0 pt}
\setlength{\parskip}{6 pt}

\onehalfspacing

\begin{document}
	
	\title{\@judultopik}
	\author{\nama \textendash \@npm} 
	
	%tulis nama dan NPM anda di sini:
	\newcommand{\nama}{Bryan Ricaldi Chandra}
	\newcommand{\@npm}{6182301013}
	\newcommand{\@judultopik}{Explorasi Reinforcement Learning menggunakan Gymnasium} % Judul/topik anda
	\newcommand{\jumpemb}{1} % Jumlah pembimbing, 1 atau 2
	\newcommand{\tanggal}{04/09/2026}
	
	% Dokumen hasil template ini harus dicetak bolak-balik !!!!
	
	\maketitle
	
	\pagenumbering{arabic}
	
	\section{Deskripsi}
	Pada tugas akhir ini, dilakukan eksplorasi, perancangan, dan implementasi algoritma \textit{Reinforcement Learning} (RL) klasik secara mandiri untuk menyelesaikan berbagai masalah pengambilan keputusan. Eksperimen ini memanfaatkan \textit{library} Python dari Farama Foundation yaitu Gymnasium API, yang menyediakan lingkungan konsisten seperti \textit{CartPole}, \textit{Taxi}, \textit{Blackjack}, serta lingkungan simulasi lainnya. Melalui pengujian pada lingkungan tersebut, performa serta perilaku agen akan dianalisis dan dibandingkan.
	
	Bidang \textit{Reinforcement Learning} (RL) terus berkembang dan memiliki minat yang kuat . Hal ini dibuktikan dengan keberhasilan pendekatan \textit{Deep Reinforcement Learning} (DRL) yang telah melebihi peforma manusia dalam berbagai permainan kompleks, seperti Go, DoTA 2, dan StarCraft 2. Namun riset terhadap RL sendiri sering kali terhambat oleh kurangnya standarisasi antarmuka, sehingga sangat sulit untuk peneliti membangun dan membandingkan model secara objektif.
	
	Gymnasium menyelesaikan permasalahan sebelumnya yaitu dengan cara memberikan pengguna antarmuka terstandardisasi yang memungkinkan penggunaan yang luas antara lingkungan simulasi dan algoritma pelatihan.  Gymnasium menyederhanakan proses pelatihan model dengan menyediakan lingkungan yang konsisten, koleksi lingkungan bawaan yang mudah dikonfigurasi, serta perangkat pendukung untuk membangun lingkungan kustom. Melalui pendekatan terpadu ini, Gymnasium secara signifikan membantu memastikan konsitensi dan ketahanan riset RL.
	
	Adapun algoritma klasik RL yang akan diimplementasikan pada tugas akhir ini adalah \textit{Q-Learning} dan \textit{SARSA}. Pemilihan kedua algoritma ini didasari oleh perbedaan fundamental dalam cara keduanya mempelajari kebijakan (\textit{policy}). \textit{Q-Learning} merupakan algoritma \textit{off-policy} yang mencari rute paling absolut dan optimal, namun cenderung berisiko tinggi dan berfluktuasi selama masa pelatihan. Sebaliknya, \textit{SARSA} adalah algoritma \textit{on-policy} yang\textbf{} memperbarui nilainya berdasarkan tindakan aktual, sehingga lebih adaptif terhadap risiko, lebih aman, dan menghasilkan grafik performa yang lebih stabil. Untuk mengevaluasi keandalan dari kedua algoritma tersebut, agen cerdas akan diuji pada beberapa lingkungan simulasi bawaan Gymnasium yang memiliki tingkat kompleksitas yang berbeda-beda.
	
	% Satu paragraf ngejelasin bedanya algoritma Sarsa dan QLearning menurut buku sunttto 
	% Satu paragraf menjelaskan tentang Markov Decision Model
	
	\section{Rumusan Masalah}
	\begin{enumerate}
		\item Bagaimana implementasi algoritma RL klasik secara mandiri agar dapat berinteraksi dengan Gymnasium API?
		\item Bagaimana perbandingan agen menggunakan algoritma klasik seperti Q-Learning dan SARSA  ketika diuji pada lingkungan yang berbeda beda?
		\item Bagaimana pengaruh \textit{hyperparameter} seperti learning rate dapat mempengaruhi kecepatan belajar dan keputusan agen?
	\end{enumerate}
	
	\section{Tujuan}
	\begin{enumerate}
		\item Mengimplementasikan algoritma RL klasik secara mandiri agar dapat beiinteraksi dengan Gymnasium API.
		\item Menganalisis dan membadingkan peforma agen yang dibangun dengan algoritma klasik seperti Q-Learning dan SARSA ketika diuji pada lingkungan yang berbeda beda.
		\item Mengevaluasi pengaruh \textit{hyperparameter} seperti \textit{learning rate} dapat mempengaruhi kecepatan belajar dan keputusan agen.
	\end{enumerate}
	
	\section{Deskripsi Perangkat Lunak}
	Tuliskan deksripsi dari perangkat lunak yang akan anda hasilkan. Apa saja fitur yang disediakan oleh PL tersebut dan apa saja kemampuan dari PL tersebut. Perhatikan contoh di bawah ini:
	
	Perangkat lunak akhir yang akan dibuat memiliki fitur minimal sebagai berikut:
	\begin{itemize}
		\item Pengguna dapat melihat denah Musem Geologi Bandung dalam bidang dua dimensi. Sedangkan pengunjung direpresentasikan menggunakan lingkaran-lingkaran kecil (tidak menggunakan gambar manusia yang diambil dari atas)
		\item Pengguna dapat memunculkan atau menghilangkan gambar {\it flow tiles} pada denah museum. 
		\item Pengguna dapat mengatur jalannya simulasi: memulai(start) simulasi, menunda(pause) simulasi, melanjutkan(continue) simulasi, maupun menghentikan(stop) simulasi
		\item Pengguna dapat mengatur banyaknya pengunjung di dalam museum, baik melalui perubahan frekuensi kedatangan pengunjung maupun menambahkan dan menghapus pengunjung satu-persatu secara manual.
		\item Posisi kamera dapat diubah (pergerakan di bidang tiga dimensi) sehingga pengguna dapat melihat simulasi di museum dari berbagai arah. 
		\item Posisi kamera dapat diubah untuk emngikuti perjalanan seorang pengunjung di dalam 
		\item Pengguna dapat memilih apakah akan menggunakan teknik {\it flow tiles} atau tidak pada saat simulasi berlangsung
		\item Jenis {\it flow tiles} yang digunakan dapat diubah-ubah pada saat simulasi sedang berlangsung
		
	\end{itemize}
	
	\section{Detail Pengerjaan Tugas Akhir}
	Tuliskan bagian-bagian pengerjaan skripsi secara detail. Bagian pekerjaan tersebut mencakup awal hingga akhir skripsi, termasuk di dalamnya pengerjaan dokumentasi skripsi, pengujian, survei, dll.
	
	\begin{enumerate}
		\item \textbf{Tahap Persiapan dan Studi Literatur}
		\begin{itemize}
			\item Melakukan studi literatur terkait konsep dasar algoritma dasar RL
			\item Mempelajari dokumentasi library Gymnasium, khususnya mengenai struktur API
			\item Mempelajari bahasa Python dan semua library yang dibutuhkan seperti Numpy
		\end{itemize}
		
		
	\end{enumerate}
	
\section{Rencana Kerja}
Rincian capaian yang direncanakan pada penyelesaian Tugas Akhir 1 (Bab 1 hingga Bab 3) adalah sebagai berikut:
\begin{enumerate}
	\item Melakukan studi literatur secara mendalam mengenai konsep RL
	\item Melakukan studi mandiri terhadap bahasa pemrograman Python
	\item Mempelajari dokumentasi dan kerangka kerja dari Gymnasium API
	\item Merumuskan parameter pengujian (\textit{hyperparameter}) seperti \textit{learning rate}, \textit{discount factor}, dan probabilitas eksplorasi ($\epsilon$), serta menentukan metrik evaluasi performa (seperti imbalan kumulatif dan konvergensi).
	\item Menyusun dokumen laporan Tugas Akhir 1
\end{enumerate}
	
	Sedangkan yang akan diselesaikan di Tugas Akhir 2 adalah sebagai berikut:
	\begin{enumerate}
		\item Mengimplementasikan algoritma RL, yaitu \textit{Q-Learning} dan \textit{SARSA}
		\item Mengimplementasikan algoritma yang telah dibuat agar dapat berinteraksi langsung dengan berbagai lingkungan simulasi dari library Gymnasium.
		\item Melakukan eksperimen dan pengujian performa agen menggunakan berbagai variasi \textit{hyperparameter}.
		\item Menganalisis hasil eksperimen secara komparatif 
		\item Menyusun dan finalisasi dokumen laporan akhir Tugas Akhir.
	\end{enumerate}
	
	\vspace{1cm}
	\centering Bandung, \tanggal\\
	\vspace{2cm} \nama \\ 
	\vspace{1cm}
	
	Menyetujui, \\
	\ifdefstring{\jumpemb}{2}{
		\vspace{1.5cm}
		\begin{centering} Menyetujui,\\ \end{centering} \vspace{0.75cm}
		\begin{minipage}[b]{0.45\linewidth}
			% \centering Bandung, \makebox[0.5cm]{\hrulefill}/\makebox[0.5cm]{\hrulefill}/2013 \\
			\vspace{2cm} Nama: \makebox[3cm]{\hrulefill}\\ Pembimbing Utama
		\end{minipage} \hspace{0.5cm}
		\begin{minipage}[b]{0.45\linewidth}
			% \centering Bandung, \makebox[0.5cm]{\hrulefill}/\makebox[0.5cm]{\hrulefill}/2013\\
			\vspace{2cm} Nama: \makebox[3cm]{\hrulefill}\\ Pembimbing Pendamping
		\end{minipage}
		\vspace{0.5cm}
	}{
		% \centering Bandung, \makebox[0.5cm]{\hrulefill}/\makebox[0.5cm]{\hrulefill}/2013\\
		\vspace{2cm} Nama: Lionov, Ph.D\\ Pembimbing Tunggal
	}
\end{document}
